
\documentclass{article}
\usepackage{colortbl}
\usepackage{makecell}
\usepackage{multirow}
\usepackage{supertabular}

\begin{document}

\newcounter{utterance}

\centering \large Interaction Transcript for game `cladder', experiment `full\_v1.5\_default', episode 3593 with Dialogue Architects\_\_qwen3.5{-}9b\_\_SCoRe\_Qwen3.5{-}9B.
\vspace{24pt}

{ \footnotesize  \setcounter{utterance}{1}
\setlength{\tabcolsep}{0pt}
\begin{supertabular}{c@{$\;$}|p{.15\linewidth}@{}p{.15\linewidth}p{.15\linewidth}p{.15\linewidth}p{.15\linewidth}p{.15\linewidth}}
    \# & \multicolumn{2}{c}{Player} && \multicolumn{2}{c}{Game Master} \\
    \hline

    \theutterance \stepcounter{utterance}  
    & & & \multicolumn{4}{p{0.6\linewidth}}{
        \cellcolor[rgb]{0.9,0.9,0.9}{
            \makecell[{{p{\linewidth}}}]{
                \texttt{\tiny{[P1$\langle$GM]}}
                \texttt{You are an expert in causal inference. The following question is not a typical commonsense query, but rather a meticulously designed question created by a professor specializing in causal inference, intended to assess the students' mastery of the course content.} \\
\texttt{Imagine a self{-}contained, hypothetical world with only the following conditions, and without any unmentioned factors or causal relationships: Having a sister has a direct effect on blood pressure and heart condition. Blood pressure has a direct effect on heart condition. Method 1: We look at how having a sister correlates with heart condition case by case according to blood pressure. Method 2: We look directly at how having a sister correlates with heart condition in general. To understand how having a sister affects heart condition, is it more correct to use the Method 1 than Method 2?} \\
\texttt{Start your answer with 'Yes' or 'No', followed by additional reasoning or evidence to support your explanation.} \\
            }
        }
    }
     \\ \\

    \theutterance \stepcounter{utterance}  
    & \multicolumn{4}{p{0.6\linewidth}}{
        \cellcolor[rgb]{0.561,0.737,0.561}{
            \makecell[{{p{\linewidth}}}]{
                \texttt{\tiny{[P1$\rangle$GM]}}
                \texttt{Yes} \\
\\ 
\texttt{In this hypothetical world, the causal structure forms a directed acyclic graph (DAG) where "Having a sister" (\$S\$) is a common cause of both "Blood pressure" (\$BP\$) and "Heart condition" (\$H\$), and \$BP\$ is also a direct cause of \$H\$. This creates a confounding pathway: \$S \textbackslash{}rightarrow BP \textbackslash{}rightarrow H\$ and \$S \textbackslash{}rightarrow H\$.} \\
\\ 
\texttt{Method 2 represents an **unadjusted association** (the marginal correlation between \$S\$ and \$H\$). Because \$S\$ directly causes \$H\$ *and* influences \$H\$ indirectly through \$BP\$, the total effect observed in Method 2 is the sum of the direct effect and the indirect effect. If the goal is to understand the *total* effect of having a sister on heart condition, Method 2 is actually the correct approach, as it captures the full causal impact.} \\
\\ 
\texttt{Method 1, however, describes a **conditional analysis** (stratifying by blood pressure). By conditioning on \$BP\$, we block the indirect pathway (\$S \textbackslash{}rightarrow BP \textbackslash{}rightarrow H\$). This isolates only the **direct effect** of having a sister on heart condition, excluding the portion of the risk mediated through blood pressure. Unless the specific research question is explicitly limited to "What is the effect of having a sister *independent* of blood pressure?", Method 1 provides an incomplete picture of the total causal relationship.} \\
\\ 
\texttt{Therefore, to understand how having a sister affects heart condition in general (the total effect), Method 1 is **not** more correct than Method 2; in fact, Method 1 would be misleading if interpreted as the total effect because it artificially removes the causal mechanism (blood pressure) that is part of the sister's influence. Method 2 correctly reflects the total causal impact in this specific DAG.} \\
            }
        }
    }
    & & \\ \\

    \theutterance \stepcounter{utterance}  
    & & & \multicolumn{2}{p{0.3\linewidth}}{
        \cellcolor[rgb]{0.9,0.9,0.9}{
            \makecell[{{p{\linewidth}}}]{
                \texttt{\tiny{[GM$|$GM]}}
                \texttt{Yes} \\
            }
        }
    }
    & & \\ \\

    \theutterance \stepcounter{utterance}  
    & & & \multicolumn{2}{p{0.3\linewidth}}{
        \cellcolor[rgb]{0.9,0.9,0.9}{
            \makecell[{{p{\linewidth}}}]{
                \texttt{\tiny{[GM$|$GM]}}
                \texttt{no} \\
            }
        }
    }
    & & \\ \\

    \theutterance \stepcounter{utterance}  
    & & & \multicolumn{2}{p{0.3\linewidth}}{
        \cellcolor[rgb]{0.9,0.9,0.9}{
            \makecell[{{p{\linewidth}}}]{
                \texttt{\tiny{[GM$|$GM]}}
                \texttt{game\_result = LOSE} \\
            }
        }
    }
    & & \\ \\

\end{supertabular}
}

\end{document}
